% The flux-anchored equilibrium loop: the condition, Step A -> Step B, and the
% convergence test that closes it.
%
% Layout note (2026-07): the previous version put the "repeat" label directly
% on the coral feedback line, and sprouted the "check each cycle" arrow out of
% that same line rather than out of Step B. Both are fixed here: the test now
% hangs below Step B on its own vertical, and the feedback runs underneath
% everything with its label in clear space above the line.
\begin{tikzpicture}[
  font=\footnotesize,
  every node/.style={align=left},
  bx/.style={rectangle, draw=black, line width=0.6pt, fill=white,
             inner sep=5pt, text width=58mm, minimum height=20mm,
             anchor=north west},
  cnd/.style={rectangle, draw=coral, line width=0.9pt, fill=coral!7,
              inner sep=5pt, text width=142mm, anchor=north west},
  result/.style={rectangle, draw=forest, line width=0.7pt, fill=forest!6,
                 inner sep=4pt, text width=62mm, anchor=north},
  arr/.style={-{Stealth[length=2mm]}, line width=0.6pt},
  loop/.style={-{Stealth[length=2mm]}, line width=0.8pt, color=coral,
               rounded corners=2pt}
]
% --- the precondition, stated before the loop it governs ------------------
\node[cnd] (C) at (0,0) {%
  {\color{coral}\textbf{Condition}} --- the loop is legal only in
  \emph{periodic steady state}: averaged over one lunation
  $\langle\partial_t T\rangle=0$, so the cycle-mean flux
  $\langle K\,\partial_z T\rangle=Q_b$ at \emph{every} depth.
  The anchor $z_0=0.55$ m must sit below the rectification zone, where
  $u_{\rm rect}<1\%$ of $Q_b$.};

\node[bx, below=7mm of C.south west, anchor=north west] (A)
  {\textbf{Step A --- forward solve}\\[2pt]
   {\color{dim}Consumes:} $T_{\rm init}(z)$\\
   {\color{dim}Action:} run $n_{\rm in}$ lunations on the \emph{skin} (CN+Newton+Thomas)\\
   {\color{dim}Produces:} $T_{\rm mean}(\text{skin}),\;T_{\rm mean}(z_0),\;u_{\rm rect}$};
\node[bx, right=26mm of A] (B)
  {\textbf{Step B --- closure reconstruction}\\[2pt]
   {\color{dim}Consumes:} $T_{\rm mean}(z_0),\;u_{\rm rect}$\\
   {\color{dim}Action:} integrate
   $\frac{d\langle T\rangle}{dz}=\frac{Q_b-u_{\rm rect}}{K(\langle T\rangle,z)}$
   downward (RK2)\\
   {\color{dim}Produces:} $T_{\rm recon}(z)\;\;(z>z_0)$};

% forward: A -> B, carrying the anchor temperature and the eddy flux
\draw[arr] (A.east) -- (B.west)
  node[midway, above, font=\scriptsize\itshape, color=dim]
  {$T_{\rm mean}(z_0),\,u_{\rm rect}$};

% the convergence test hangs off Step B on its own vertical, because that is
% what it tests -- it is not a branch of the feedback line
\node[result, below=13mm of B.south] (R)
  {\textbf{Converged?}\quad
   $\bigl|T_{\rm mean}(z_0)^{(k)}-T_{\rm mean}(z_0)^{(k-1)}\bigr|<0.005$\,K\\[1pt]
   {\color{dim}typically 4 outer cycles: 2 primer $+$ 2 production}};
\draw[arr] (B.south) -- (R.north)
  node[midway, right=1mm, font=\scriptsize\itshape, color=dim] {check};

% feedback: not yet converged -> go round again. Three explicit segments so the
% arrow rises VERTICALLY into Step A's bottom edge and lands flush on it. The
% earlier `-|' form finished with a horizontal hop, which drove the head along
% the border and over the box text.
\coordinate (lo)  at ($(R.south)+(0,-9mm)$);
\coordinate (loA) at (A.south |- lo);
\draw[loop] (R.south) -- (lo) -- (loA) -- (A.south);
\node[font=\scriptsize\itshape, color=coral, anchor=south, inner sep=1pt,
      fill=white]
  at ($(lo)!0.5!(loA)+(0,1.2mm)$) {no --- set $T_{\rm init}=T_{\rm recon}$, repeat};
\end{tikzpicture}
